\scanpage{319}
\mspage{235}
% archive sheet number: 318

On voit, par les $=$ arguments, que $I$ \deleted{\unsure{est}}satisfait $\mathrm{PF}1'$ i.e. est ps. \unsure{filtr.te|filtrante}. Ainsi, si on introduit la condition

\tikz[baseline=(b.base)]\node[draw,ellipse,inner sep=2pt](b){$b''$}; $I'$ est filtrante, \deleted{\unsure{$I$ et}} les $ {}_{i'}\backslash I$ $(i'\in\mathrm{Ob}\,I')$ sont pseudofiltrantes \inserted{(i.e. satisfont $\mathrm{PF}1'$)} (\deleted{donc} $f$ cofinal).

On trouve les implications

\[(a\Leftrightarrow b\Leftrightarrow b')\Rightarrow(b''\Leftrightarrow c)\]

et \deleted{\unsure{que}} les conditions équivalentes c), b'') impliquent que $I$ est \inserted{également} ps. filtrante, et que pour \unsure{tt|toute|$\forall$} double-flèche \inserted{\diagram{\begin{tikzcd}[ampersand replacement=\&,cramped,column sep=small] \unsure{\zeta|\xi}_0\arrow[r,shift left=.6ex,"u_1"]\arrow[r,shift right=.6ex,"u_2"']\&\zeta_1 \end{tikzcd}}} dans un $ {}_{i'}\backslash I$ ou dans $I$, il existe une flèche $\zeta_1\overset{v}{\longrightarrow}\zeta_2$ dans la $=$ catégorie, telle que $vu_1$ et $vu_2:$ \diagram{\begin{tikzcd}[ampersand replacement=\&,cramped,column sep=small] \zeta_0\arrow[r,shift left=.6ex]\arrow[r,shift right=.6ex]\&\zeta_2 \end{tikzcd}} aient \uline{$=$ image} \uline{dans $I'$}.

Donc si $f$ est fidèle, alors ces conditions \unsure{a)|c)} c' impliquent que $I$ et les $ {}_{i'}\backslash I$ satisfont aussi la condition $\mathrm{PF}2$ de la double flèche, donc sont filtrantes, \deleted{donc} dans ce cas les cinq conditions a) b) b'), c) c') sont équivalentes. Mais aussi lorsque $I$ est ordonnée, \deleted{\unsure{donc|de}} \unsure{sorte|suite|sûr} que la condition de la double flèche dans $I$ est satisfaite automatiquement.

